% 第2章: モジュールという考え方
\section{モジュールという考え方}

%% ── Q1: モジュールとは ──────────────────────
\begin{qabox}{モジュールって何？}

モジュールとは、「1つのハードウェア部品の制御を、1つのクラスにまとめたもの」です。

ルールはシンプルです。

\begin{center}
\begin{tcolorbox}[
  enhanced, colback=accentblue!8, colframe=accentblue,
  arc=2mm, width=10cm, breakable,
]
  \centering\large\bfseries
  1ハードウェア = 1クラス（1モジュール）
\end{tcolorbox}
\end{center}

第1章のベタ書きコードから、LED部分だけを抜き出してクラスにしてみましょう。

\begin{goodexample}{LedModuleを抽出}
\begin{lstlisting}
// LedModule.h
class LedModule {
private:
    uint8_t _pin;
    bool _state;

public:
    LedModule(uint8_t pin) : _pin(pin), _state(false) {}

    bool init() {
        pinMode(_pin, OUTPUT);
        return true;
    }

    void turnOn() {
        _state = true;
        digitalWrite(_pin, HIGH);
    }

    void turnOff() {
        _state = false;
        digitalWrite(_pin, LOW);
    }

    void toggle() {
        _state ? turnOff() : turnOn();
    }
};
\end{lstlisting}
\end{goodexample}

これで main.cpp は次のようにすっきりします。

\begin{lstlisting}
LedModule led1(2);   // ピン2のLED
LedModule led2(4);   // ピン4のLED

void setup() {
    led1.init();
    led2.init();
}
\end{lstlisting}

LEDの制御に関するコードは \texttt{LedModule} の中に閉じ込められました。main.cppを見ても「LEDが2つあって、initしている」ということが一目で分かります。

\end{qabox}

%% ── Q2: Config構造体 ──────────────────────
\begin{qabox}{なんでConfig構造体が必要なの？}

先ほどの例では、コンストラクタにピン番号だけ渡しました。でも設定項目が増えるとどうなるでしょう？

\begin{badexample}{引数が増えすぎる}
\begin{lstlisting}
// ピン番号、初期状態、PWMチャンネル、PWM周波数...
LedModule led(2, false, 0, 5000);
// どの引数が何を意味するか分からない！
\end{lstlisting}
\end{badexample}

\begin{goodexample}{Config構造体にまとめる}
\begin{lstlisting}
// 設定をまとめた構造体
struct LedConfig {
    uint8_t pin;
};

class LedModule {
private:
    LedConfig _config;  // 設定を丸ごと保持

public:
    LedModule(const LedConfig& config) : _config(config) {}

    bool init() {
        pinMode(_config.pin, OUTPUT);
        return true;
    }
};

// 使う側 — 何の設定か一目瞭然
const LedConfig LED1_CONFIG = { .pin = 2 };
LedModule led1(LED1_CONFIG);
\end{lstlisting}
\end{goodexample}

Config構造体を使うメリットは3つです。

\begin{enumerate}
  \item 設定項目が増えても、構造体にフィールドを追加するだけ
  \item 設定値を1か所（\texttt{ProjectConfig.h}）にまとめて管理できる
  \item コンストラクタの引数の意味が明確
\end{enumerate}

\end{qabox}

%% ── Q3: ConfigとDataの違い ──────────────────
\begin{qabox}{ConfigとDataの違いは？}

モジュールが扱う情報は、大きく2種類に分かれます。

\begin{center}
\begin{tabularx}{\textwidth}{l X X}
\toprule
& \textbf{Config} & \textbf{Data} \\
\midrule
意味 & 起動時の設定 & 実行中に変化するデータ \\
いつ決まる？ & プログラム開始前（コンパイル時） & プログラム実行中 \\
変わる？ & 基本的に変わらない & 常に変わる \\
例 & ピン番号、I2Cアドレス & LED状態、センサー値 \\
\bottomrule
\end{tabularx}
\end{center}

\begin{lstlisting}
// Config — 起動時に決まる設定
struct LedConfig {
    uint8_t pin;       // ピン番号（変わらない）
};

// Data — 実行中に変わるデータ
struct LedData {
    bool state = false;        // ON/OFF状態
    uint8_t brightness = 0;    // 明るさ
};
\end{lstlisting}

\begin{notebox}[なぜ分けるの？]
ConfigとDataを分けることで、「この値は実行中に変わるのか？」が型レベルで明確になります。Configは \texttt{const} にできるので、うっかり書き換えてしまうバグも防げます。
\end{notebox}

\end{qabox}
